\begin{center}
{{\bf\fontsize{14pt}{14.5pt}\selectfont 
KẾT LUẬN}}
\end{center}
\addcontentsline{toc}{chapter}{KẾT LUẬN}

Tiểu luận đã làm rõ ranh giới toán học giữa hình học Euclid cổ điển và hình học Origami thông qua hệ thống 7 tiên đề Huzita--Hatori, với các kết quả trọng tâm:

\begin{itemize}
    \item \textbf{Giải mã giới hạn Euclid:} Thước kẻ và compa chỉ giải quyết được các phương trình có bậc mở rộng trường số dạng $2^k$, dẫn đến sự bất khả thi đối với bài toán chia ba góc và nhân đôi khối lập phương (vốn đòi hỏi giải phương trình bậc ba tối tiểu).
    \item \textbf{Bản chất toán học của Tiên đề 6:} Thao tác đưa đồng thời hai điểm chạm vào hai đường thẳng tương đương với việc dựng tiếp tuyến chung của hai đường parabol. Cơ chế này mở rộng bậc trường số lên dạng $2^a \cdot 3^b$, cho phép tìm nghiệm thực của phương trình bậc ba trực tiếp qua tương tác vật lý mà không cần công thức đại số phức tạp.
    \item \textbf{Giá trị thực nghiệm:} Quy trình gập của Abe Hisashi và Peter Messer đã chứng minh tính khả thi tuyệt đối trên giấy thực tế. Bảng số liệu đo đạc cho thấy sai số thực nghiệm hoàn toàn nằm trong ngưỡng cho phép ($\approx \SI{0.39}{\percent} - \SI{1.33}{\percent}$), xuất phát chủ yếu từ độ co giãn và độ dày của vật liệu giấy.
\end{itemize}

Vượt ra ngoài một nghệ thuật tạo hình dân gian, Origami chính là một hệ tiên đề toán học chặt chẽ. Sự giao thoa giữa những nếp gấp mộc mạc và cấu trúc đại số trừu tượng đã minh chứng rằng: giới hạn của một bài toán đôi khi không nằm ở bản chất hình học, mà chỉ nằm ở công cụ ta lựa chọn để tiếp cận nó.